% page 269 du fac-simile (image p-268 du scan)
% bloc 160.0 x 242.0 mm ; marge gauche 42.0 mm
\pgc{20.320mm}{0.000mm}{
\l{\cel{54}261.}
\l{}
\l{\sou{kaoso}. I. Konfuzeso di la univers-elementi qui, segun ula}
\l{\cel{3}teogonii, existis ante la organizeso di mondo. - II.}
\l{\cel{3}(\sou{mataf.}) Konfuzeso e desordineso absoluta. - DEFIRS.}
\l{}
\l{\sou{kapo}. (\sou{anat.}) I. Parto supra di la homo-korpo - parto avana di}
\l{\cel{3}la animalo-korpo - qua kontenas la cerebro e la organi pre-}
\l{\cel{3}cipua di la sensi. - II. (\sou{metaf}.) Parto supra : l. to quo}
\l{\cel{3}direktas ; 2. somito e (precipue) somito kalotatra; 3. (me-}
\l{\cel{3}\sou{taf.}) parto avana : to quo esas situita an-avan la objekto ;}
\l{\cel{3}4. to quo komencas ulo (kom ex.: la kapo di relvoyo-lineo).}
\l{\cel{3}- ISL.}
\l{}
\l{\sou{kapabla}. (La persono) qua povas (mentale o korpale) exekutar}
\l{\cel{3}bone to quon onu konfidas a lu kom agenda o facenda.-DEFIS.}
\l{}
\l{\sou{kapaca}. (\sou{extense}\cel{1}di). Tale granda ke lu povas kontenar (ta o}
\l{\cel{3}ca quanto de ulo). - DEFIS.}
\l{}
\l{kapear.\cel{2}(\sou{netrans.})\cel{2}Navigar (dum vento forta) transverse en}
\l{\cel{3}vento, sub seglaro mikra, la segli bursatrigite. - FIS.}
\l{}
\l{\sou{kapelo}. (\sou{kristanismo}). Loko sakrigita, ube onu konservas la}
\l{\cel{3}mantelo, la restaji de santo. - II. l. Loko konsakrita a la}
\l{\cel{3}kulto, en kastelo, kongregaciono, hospico, kolegio, e c. ;}
\l{\cel{3}la personi qui oficias en ta loko.\cel{2}2. Parto di kirko qua}
\l{\cel{3}kontenas altaro partikulara exter la koreyo. 3. Kirko qua}
\l{\cel{3}ne havas ankore la titulo \textquotedbl{}parokiala\textquotedbl{}. - DEFIS.}
\l{}
\l{kapelino. Kofio por muliero, qua decensas til an-sur elua}
\l{\cel{3}shultri. - FIRS.}
\l{}
\l{\sou{kapero.}\cel{2}(\sou{bot.)}\cel{2}Yuna burjoneto de arboreto (L.\cel{1}\sou{capparia spi}}
\l{\cel{3}nosa)\cel{2}kultivata sude di Europa, qua, konfitite en vinagro,}
\l{\cel{3}uzesas kom kondimento. - DEFIRS.}
\l{}
\l{\sou{kapilara}. Di qua la diametro esas haro-diametra. - DEFIRS.}
\l{}
\l{kapilario. (\sou{bot.)}\cel{2}Filiko di qua la stipito e la foliaro esas}
\l{\cel{3}tenua ; e di qua la folii uzesas sive infuze, sive sirope.}
\l{\cel{3}- L.\cel{1}\sou{adiantus capillus veneris}.}
\l{}
\l{kapistro.\cel{2}(\sou{kirurg.}) Bandajo per qua on reduktas la rupturi}
\l{\cel{3}e la luxacuri che la osti di la mandibulo. - DFIS.}
\l{}
\l{kapitalo. (\sou{financo}). La pekunio quan ulu posedas (kontraste kun}
\l{\cel{3}la revenuo, qua esas acesora). La pekunio,egardata kom pro-}
\l{\cel{3}dukto-moyeno. - DEFIRS.}
\l{}
\l{kapitano. I. La persono qua komandas en regimento-kompanio.}
\l{\cel{3}II. La persono qua komandas sur milito-navo, kom chefo,}
\l{\cel{3}oficiro, suprega. - DEFIRS.}
\l{}
\l{kapitelo. (\sou{arkitekt.})Parto supra di kolono, di pilastro, qua}
\l{\cel{3}esas quaze la krono di la fusto. - DEFIRS.}
}
